\documentclass[../../main.tex]{subfiles}

\begin{document}

\section{Vectors}

We have discussed on vectors more in-depth in the part for Linear
Algebra. In this section however, let's focus more on the geometric
intuition than the abstract nature of vectors.

\begin{definition}{}{}
\textbf{Vectors} are quantities with both magnitude and a direction.
\end{definition}

Some example of vectors include force, velocity, and acceleration
because they all have a specific direction and also a magnitude.
When we graph them, we can think of them as a arrow with starting
and ending points.

\begin{figure}[H]
\centering
\tdplotsetmaincoords{70}{110}
\begin{tikzpicture}[
        tdplot_main_coords,
        axis/.style={->,thick},
        scale=0.4
    ]
	\draw[axis] (-8,0,0) -- (8,0,0) node[left] {$x$};
	\draw[axis] (0,-7,0) -- (0,8,0) node[right] {$y$};
    \draw[axis] (0,0,-5) -- (0,0,8) node[left] {$z$};

    \draw[axis] (2,-3,2) -- (1,3,5) node[above right] {$\mathbf{v}$};
    \draw[axis] (2,-3,-1) -- (1,3,2) node[above right] {$\mathbf{w}$};
    \node[circle,fill=black,inner sep=0pt,minimum size=1mm]
        at (2,-3,2) {};
    \node[circle,fill=black,inner sep=0pt,minimum size=1mm]
        at (2,-3,-1) {};
\end{tikzpicture}
\end{figure}

From now, let's represent vectors with $\mathbf{u}, \mathbf{v},
\mathbf{w}$ and scalars, or real numbers, with $a, b, \alpha, \beta$.
From the graph above, we can define another definition.

\begin{definition}{}{}
The \textbf{terminal point} of a vector is the point where the vector
ends. The \textbf{initial point} is the point where the vector begins.
\end{definition}

When we want to add two vectors $\mathbf{v}$ and $\mathbf{w}$,
we would connect the terminal point of $\mathbf{v}$ by the starting
point of $\mathbf{w}$ and obtain a new vector by connecting the
initial point of $\mathbf{v}$ and the terminal point of $\mathbf{w}$
as shown by the figure below.

\begin{figure}[H]
\centering
\tdplotsetmaincoords{70}{110}
\begin{tikzpicture}[
        tdplot_main_coords,
        axis/.style={->,thick},
        scale=0.4
    ]
	\draw[axis] (-8,0,0) -- (8,0,0) node[left] {$x$};
	\draw[axis] (0,-7,0) -- (0,8,0) node[right] {$y$};
    \draw[axis] (0,0,-5) -- (0,0,8) node[left] {$z$};

    \draw[axis] (2,-3,-1) -- (1,-1,2) node[above] {$\mathbf{v}$};
    \draw[axis] (1,-1,2) -- (1,6,4) node[above] {$\mathbf{w}$};
    \draw[axis, Red] (2,-3,-1) -- (1,6,4)
        node[below right] {$\mathbf{v + w}$};
    \node[circle,fill=black,inner sep=0pt,minimum size=1mm]
        at (2,-3,-1) {};
    \node[circle,fill=black,inner sep=0pt,minimum size=1mm]
        at (1,-1,2) {};
\end{tikzpicture}
\end{figure}

Multiplying a vector by scalar $k$ can be though of as lengthening
the vector by a factor of $k$. If $k < 0$, the vector is pointing
the opposite direction.

\begin{figure}[H]
\centering
\tdplotsetmaincoords{70}{110}
\begin{tikzpicture}[
        tdplot_main_coords,
        axis/.style={->,thick},
        scale=0.4
    ]
    \draw[axis] (-8,0,0) -- (8,0,0) node[left] {$x$};
    \draw[axis] (0,-7,0) -- (0,8,0) node[right] {$y$};
    \draw[axis] (0,0,-5) -- (0,0,8) node[left] {$z$};

    \draw[axis] (2,-3,2) -- (1.5,0,3.5) node[above] {$\mathbf{v}$};
    \draw[axis] (-2,3,-2) -- (-1.5,-0,-3.5)
        node[right] {$\mathbf{-v}$};
    \draw[axis] (2,-3,-1) -- (1,3,2) node[below] {$\mathbf{2v}$};
    \node[circle,fill=black,inner sep=0pt,minimum size=1mm]
        at (2,-3,2) {};
    \node[circle,fill=black,inner sep=0pt,minimum size=1mm]
        at (2,-3,-1) {};
\end{tikzpicture}
\end{figure}

Also for the purpose of the notes, let's consider the vectors
with same magnitude and direction but different initial vectors
as the same vectors. The being said, we can write any vectors
in $n$-dimension as $\langle v_1, v_2, v_3, \ldots, v_n \rangle$
where $(v_1, v_2, v_3, \ldots, v_n)$ is the terminal point and
$(0, 0, \linebreak, 0 \ldots, 0)$ is the initial point.
Such numbers $v_1, \ldots, v_n$ are known as the \textbf{components}
of $\mathbf{v}$.

Now that we have defined $-\mathbf{v}$, we can consider vector
subtraction. We can consider vector subtraction as adding two vectors
where one is scaled by a negative factor.

\begin{figure}[H]
\centering
\tdplotsetmaincoords{70}{110}
\begin{tikzpicture}[
        tdplot_main_coords,
        axis/.style={->,thick},
        scale=0.4
    ]
    \draw[axis] (-8,0,0) -- (8,0,0) node[left] {$x$};
    \draw[axis] (0,-7,0) -- (0,8,0) node[right] {$y$};
    \draw[axis] (0,0,-5) -- (0,0,8) node[left] {$z$};

    \draw[axis] (2,2,-2) -- (1,5,3)
        node[right] {\footnotesize $\mathbf{v}$};
    \draw[axis] (2,2,-2) -- (2,-2,2)
        node[left] {\footnotesize $\mathbf{w}$};
    \draw[axis,dotted] (1,5,3) -- (1,1,7);
    \draw[axis,dotted] (2,-2,2) -- (1,1,7);
    \draw[axis,Red] (2,-2,2) -- (1,5,3)
        node[pos=0.7, below] {\footnotesize $\mathbf{v - w}$};
    \draw[axis,Blue] (2,2,-2) -- (1,1,7)
        node[pos=0.7, left] {\footnotesize $\mathbf{v + w}$};
\end{tikzpicture}
\end{figure}

Using components, we can write addition, subtraction and scalar
multiplication as the following for $\mathbf{v} = \langle v_1, v_2,
v_3 \rangle$ and $\mathbf{w} = \langle w_1, w_2, w_3 \rangle$.
\begin{align*}
    \mathbf{v} + \mathbf{w}
    &= \langle v_1 + w_1, v_2 + w_2, v_3 + w_3 \rangle \\
    \mathbf{v} - \mathbf{w}
    &= \langle v_1 - w_1, v_2 - w_2, v_3 - w_3 \rangle \\
    k\mathbf{v}
    &= \langle kv_1, kv_2, kv_3 \rangle
\end{align*}
As you can see from the equations above, we can see that vector
addition is commutative and associative and scalar multiplication is
distributive. I will keep them brief here because we already discussed
about these in-depth on my other notes for linear algebra. Now this
property leads to the natural question of subtracting a vector by
itself.

\begin{definition}{}{}
The \textbf{zero vector} is a vector with magnitude of zero.
\end{definition}

Intuitively, we can think of a zero vector as vector with the same
initial and terminal points. Another way of thinking about a zero
vector is a vector with zero as its components. This will naturally
leads to $\mathbf{v} + \mathbf{0} = \mathbf{v}$ for any vector
$\mathbf{v}$. For the purpose of this notes, let's consider zero
vector as the vector that is parallel to every vectors.

Continuing with the definition of components and the magnitude of a
vector, we can define an important terminology.

\begin{definition}{}{}
The magnitude of a vector $\mathbf{v}$ is known as the \textbf{norm}
of $\mathbf{v}$, denoted as $\| \mathbf{v} \|$.
\end{definition}

Notice that the magnitude, or length, cannot be negative. Therefore,
we can see that the norm of any vector is non-negative. Moreover,
$\| \mathbf{0} \| = 0$. Also note that scalar multiplication is
defined element-wise. Using the distance formula, we can also obtain
the following very important equation for a vector $\mathbf{v}$ and
scalar $k$.
\[
\| k \mathbf{v} \| = |k| \cdot \| \mathbf{v} \|
\]
Let's take a look at a quick example. Consider the vector
$\mathbf{v} = \langle 1, 2, 3 \rangle$. Notice that the following
equation holds.
\[
\| \mathbf{v} \|
= \sqrt{1^2 + 2^2 + 3^2}
= \sqrt{14}
\]
We can also see that the following equation holds with scalar $2$
and $2 \langle 1, 2, 3 \rangle = \langle 2, 4, 6 \rangle$.
\begin{align*}
    \| 2 \mathbf{v} \|
    &= \sqrt{2^2 + 4^2 + 6^2}
    = \sqrt{2^2 (1^2 + 2^2 + 3^2)} \\
    &= 2 \sqrt{1^2 + 2^2 + 3^2} \\
    &= 2 \cdot \| \mathbf{v} \|
\end{align*}
With this equation in mind, let's consider the following question.

\begin{important}
Given a vector $\mathbf{v}$, how do we find a unit vector, or
vector with length $1$, with the same direction as $\mathbf{v}$?
\end{important}

One way of thinking this is rearranging our equation
$\| k \mathbf{v} \| = |k| \cdot \| \mathbf{v} \|$. Consider the
following equations with unit vector $\mathbf{u}$ and
$\| \mathbf{v} \| = k \neq 0$.
\[
\left| \frac{1}{k} \right| \cdot \| \mathbf{v} \|
= \frac{1}{k} \cdot k
= 1
= \| \mathbf{u} \|
\]
This solves it! We can find a unit vector with the same direction
$\mathbf{v}$ by dividing it by its norm. Another variation of this
equation would be the following.
\[
\| \mathbf{v} \| = |k| \cdot \| \mathbf{u} \|
\]
\begin{important}
This implies that any vector $\mathbf{v}$ can be represented as
the product of its norm and a unique unit vector with the same
direction as $\mathbf{v}$. Such process of multiplying a nonzero
vector by the reciprocal of its norm is known as \textbf{normalizing}.
\end{important}

Finally, another standard way of writing vectors is by linear
combination of the standard basis vectors defined as the following.
\[
\mathbf{i} = \langle 1, 0, 0 \rangle, \quad
\mathbf{j} = \langle 0, 1, 0 \rangle, \quad
\mathbf{k} = \langle 0, 0, 1 \rangle
\]
With this, we can write $\mathbf{v} = \langle 1, 2, 3 \rangle$ as
$\mathbf{i} + 2\mathbf{j} + 3\mathbf{k}$.

More detailed and abstract introduction of vectors are included
in the linear algebra part of LibreNotebook, and I highly recommend
checking that out! From here, we will discuss more foundational and
geometric interpretations of vectors that we have not discussed
in my notes for linear algebra.

\subsection{Dot Products and Projections}

As always, let's start with definition.

\begin{definition}{}{}
The \textbf{dot product} of two vectors $\mathbf{u}$ and $\mathbf{v}$
of the same dimension is defined as the sum of all products of
elements in the same index.
\end{definition}

Writing them with mathematical notations, we can write them as the
following.
\[
\langle u_1, u_2, \ldots, u_n \rangle \cdot
    \langle v_1, v_2, \ldots, v_n \rangle
= \sum_{i=1}^n u_i v_i
\]
Here is a quick example.
\[
\langle 1, 2, 3 \rangle \cdot \langle -1, 0, 1 \rangle
= -1 + 0 + 3
= 2
\]
Continuing, below are some properties of dot products. The statements
below are quite self-evident as they mirror the properties of addition
and subtraction.

\begin{theorem}{}{Properties of the Dot Product}
For vectors $\mathbf{u}, \mathbf{v}, \mathbf{w}$ of the same dimension
and an arbitrary scalar $\lambda$, the following equations hold.
\begin{enumerate}
\item $\mathbf{u} \cdot \mathbf{v} = \mathbf{v} \cdot \mathbf{u}$
\item $\lambda (\mathbf{u} \cdot \mathbf{v})
= (\lambda \mathbf{u}) \cdot \mathbf{v}
= \mathbf{u} \cdot (\lambda \mathbf{v})$
\item $\mathbf{u} \cdot (\mathbf{v} + \mathbf{w})
= \mathbf{u} \cdot \mathbf{v} + \mathbf{u} \cdot \mathbf{w}$
\item $\mathbf{v} \cdot \mathbf{v} = \| v \|^2$
\end{enumerate}
\end{theorem}

One of the reasons why the dot product of two vectors is so important
is because it helps us determine the angle between the two vectors.
Consider the following theorem.

\begin{theorem}{}{The Dot Product and the Angle}
Let $\mathbf{u}, \mathbf{v}$ be two nonzero vectors. Let $\theta$ be
the smaller angle between the two vectors when their initial point
coincide. Then, the following equation holds.
\[
\mathbf{u} \cdot \mathbf{v}
= \| u \| \| v \| \cos\theta
\]
\end{theorem}
\begin{proof}
First and foremost, consider the following equation by the law of
cosines.
\[
\cos\theta
= \frac{ \| \mathbf{u} \|^2 + \| \mathbf{v} \|^2 -
    \| \mathbf{v} - \mathbf{u} \|^2
}{ 2 \| \mathbf{u} \| \| \mathbf{v} \| }
\]
By Theorem \ref{theo:Properties of the Dot Product}, the following
equations hold.
\begin{align*}
    \cos\theta
    &= \frac{ \| \mathbf{u} \|^2 + \| \mathbf{v} \|^2 -
        ( \mathbf{v} - \mathbf{u} ) \cdot ( \mathbf{v} - \mathbf{u} )
    }{ 2 \| \mathbf{u} \| \| \mathbf{v} \| } \\
    &= \frac{ \| \mathbf{u} \|^2 + \| \mathbf{v} \|^2 - \left(
        \mathbf{v} \cdot \mathbf{v}  - 2 \mathbf{u} \cdot \mathbf{v}
            + \mathbf{u} \cdot \mathbf{u}
        \right)
    }{ 2 \| \mathbf{u} \| \| \mathbf{v} \| } \\
    &= \frac{2 \mathbf{u} \cdot \mathbf{v}}
        { 2 \| \mathbf{u} \| \| \mathbf{v} \| } \\
    \therefore 
    \| u \| \| v \| \cos\theta
    &= \mathbf{u} \cdot \mathbf{v}
\end{align*}
Thus, the theorem holds.
\end{proof}

Now this theorem tells us the relationship between the angle between
two vectors and their dot product.

\begin{important}
If $0 \leq \theta < \frac{\pi}{2}$, then $\mathbf{u} \cdot \mathbf{v}
> 0$. If $\theta = \frac{\pi}{2}$, then $\mathbf{u} \cdot \mathbf{v}
= 0$. If $\frac{\pi}{2} < \theta < \pi$, then $\mathbf{u} \cdot
\mathbf{v} < 0$.
\end{important}

This somewhat simple observation with the cosine function and the
fact that $\| \mathbf{u} \|, \| \mathbf{u} \| > 0$ reveals the
relationship between the dot product of two vectors and their angles.
Continuing from this observation, we can establish a definition.

\begin{definition}{}{}
Two nonzero vectors $\mathbf{u}$ and $\mathbf{v}$ are
\textbf{orthogonal} if they form a right angle. In other words,
$\mathbf{u} \cdot \mathbf{v} = 0$.
\end{definition}

Continuing with the relationship and definition, let's discuss
more on the angle of vectors.

\begin{definition}{}{}
The \textbf{direction angles} of a nonzero vector $\mathbf{v}$ is
the angle between $\mathbf{v}$ and the standard basis vectors.
\end{definition}

To find the direction angles of $\mathbf{v} = \langle v_1, v_2, v_3
\rangle$, let's return to Theorem \ref{theo:The Dot Product and the
Angle}. Consider the following equation for the angle $\alpha$
between $\mathbf{v}$ and the standard basis vector $\mathbf{i}$.
\begin{align*}
    \mathbf{v} \cdot \mathbf{i}
    &= \| \mathbf{v} \| \| \mathbf{i} \| \cos\alpha \\
    v_1 &= \| \mathbf{v} \| \cos\alpha \\
    \therefore
    \alpha &= \cos^{-1} \left( \frac{v_1}{\| \mathbf{v} \|} \right)
\end{align*}
Following the same pattern for $\beta$ and $\gamma$, the can
establish the following equations.
\[
\cos\alpha = \frac{v_1}{\| \mathbf{v} \|}, \quad
\cos\beta = \frac{v_2}{\| \mathbf{v} \|}, \quad
\cos\gamma = \frac{v_3}{\| \mathbf{v} \|}
\]
Indeed, we have special term for the values above. The values above
are known as the \textbf{direction cosines} of $\mathbf{v}$. For
vectors of direction cosines, the can write the following equation.
\[
\langle \cos\alpha, \cos\beta, \cos\gamma \rangle
= \left\langle
    \frac{v_1}{\| \mathbf{v} \|},
    \frac{v_2}{\| \mathbf{v} \|},
    \frac{v_3}{\| \mathbf{v} \|}
\right\rangle
= \frac{\mathbf{v}}{\| \mathbf{v} \|}
\]
Continuing, let's consider the case when $\| \mathbf{v} \| = 1$.
In that case,
\[
\langle \cos\alpha, \cos\beta, \cos\gamma \rangle
= \left\langle
    \frac{v_1}{\| \mathbf{v} \|},
    \frac{v_2}{\| \mathbf{v} \|},
    \frac{v_3}{\| \mathbf{v} \|}
\right\rangle
= \mathbf{v}
\]
also hold! Finally, here is another interesting properties of the
direction cosines.
\[
\cos^2\alpha + \cos^2\beta + \cos^2\gamma
= \left( \frac{v_1}{\| \mathbf{v} \|} \right)^2
    + \left( \frac{v_2}{\| \mathbf{v} \|} \right)^2
    + \left( \frac{v_3}{\| \mathbf{v} \|} \right)^2
= \frac{\| \mathbf{v} \|^2}{\| \mathbf{v} \|^2}
= 1
\]
We can see that for nonzero $\mathbf{v}$, we can generalize this
for any $n$. Therefore the sum of the squares of direction cosines
are always $1$.





\end{document}


